\documentclass[10pt,a4paper]{amsart}
\usepackage[margin=19mm]{geometry}
\usepackage{amsmath,amssymb,mathtools}
\usepackage[scr=rsfs,cal=euler]{mathalfa}
\usepackage{xfrac}
\usepackage[unicode,hidelinks,bookmarks=false]{hyperref}
\newcommand{\ie}{i.e.\ }
\newcommand{\cf}{cf.\ }
\newcommand{\resp}{respectively\ }
\newcommand{\Subsection}{\subsection}
\providecommand{\nobreakdash}{\nobreak\hbox{-}}
\setlength{\emergencystretch}{2em}
\multlinegap0pt
\usepackage{bm}
\usepackage{bbm}
\usepackage{dsfont}
\usepackage{xspace}
\usepackage{cancel}
\usepackage{mathtools}
\usepackage{cleveref}
\usepackage{amsthm}
\makeatletter
\@ifundefined{theorem}{\newtheorem{theorem}{Theorem}}{}
\makeatother
\newcommand*\diff{\mathop{}\!\mathrm{d}}
\title[Estimates for pseudo-differential operators on the torus rev]{Estimates for pseudo-differential operators on the torus revisited. I}
\author{Duv\'an Cardona, Manuel Alejandro Mart\'inez}
\date{Source: arXiv:2502.20575v1 (CC BY 4.0)}
\begin{document}
\maketitle

\noindent\emph{Source and licence.} Estimates for pseudo-differential operators on the torus revisited. I. Version arXiv:2502.20575v1 (CC BY 4.0), distributed under the Creative Commons Attribution 4.0 International licence, as recorded in the arXiv licence metadata for this exact version and shown on the abstract page https://arxiv.org/abs/2502.20575v1. Full rights notice: licenses/arxiv-2502.20575-CC-BY-4.0.txt.\par\smallskip

\noindent\textit{Editorial scope.} Counted unit: the Theorem whose source label is \texttt{theo:sigma-kernel-estimate} in the article (source0.tex lines 432--452, source label \texttt{theo:sigma-kernel-estimate}), delivered together with the article's own complete original proof of it (source0.tex lines 455--605, one \texttt{proof} environment of 151 lines). No mathematics is written, repaired or re-derived: statement and proof are the article's own text line for line. Required context retained uncounted: theo:22boundR (lines 327--335); eq:teo3-1a (lines 345--371). Every definition, display and label of those lines is retained unchanged. Omitted from this package and not redistributed: 1--326, 336--344, 372--431, 453--454, 606--1101. Nothing inside a retained excerpt is omitted or altered: every retained body is delivered exactly as the article's lines are written, so no omission receipt of any kind accompanies this package. Citations: the retained excerpts carry no citation command at all, so no bibliography entry of the article is transcribed and no reference is redistributed. Reference adaptation: none was needed and none was made; every reference inside the delivered unit resolves inside it, so no unresolved reference or citation is produced. Printed slips kept literal: no printed word, symbol, hypothesis or number of the counted statement or of its proof was altered. Layout and class: the article's own class is \texttt{amsart} with packages including amsmath, amsthm, amscd, amsfonts, amssymb, graphicx, color, mathrsfs, hyperref, slashed, cleveref; the wrapper is the lane's pinned \texttt{amsart} editorial wrapper at 10pt on A4, which restates the theorem-like environments the retained text needs and adds the article's own macro definitions that the retained text needs (\texttt{\textbackslash diff}), with extra wrapper packages (bm, bbm, dsfont, xspace, cancel, mathtools, cleveref). The delivered numbering is the unit's own. Assets: no retained span contains a figure environment, a graphics-inclusion command, a PDF-page inclusion, a table or a TikZ picture; no image file is copied into this package, the package declares no asset dependency and no image file is redistributed with it. Licence: version arXiv:2502.20575v1 of this article is distributed under the Creative Commons Attribution 4.0 International licence, as recorded in the arXiv licence metadata for this exact version and shown on the abstract page https://arxiv.org/abs/2502.20575v1. A full rights notice is delivered at licenses/arxiv-2502.20575-CC-BY-4.0.txt. Characters: every retained excerpt is pure ASCII, so the delivered engine, XeLaTeX with its default Latin Modern text font, reports no missing character.
\begin{theorem}
    Let $\tilde{p}:\mathbb{R}^n \times \mathbb{R}^n \rightarrow \mathbb{C}$ be a symbol such that for $0<\rho\leq1$, $0\leq \delta < 1$, $m\leq -n \lambda = -n\max\{0, (\delta - \rho)/2\}$ and $|\alpha|,|\beta| \leq \lceil n/2\rceil$ satisfies: 
    \begin{equation}
        \left|\partial^\alpha_\xi \partial^\beta_x \tilde{p}(x, \xi)\right| \leq C_{\alpha\beta}\langle\xi\rangle^{m-\rho|\alpha| + \delta|\beta|}
    \end{equation}

    Then $\mathrm{Op}(\tilde{p})$ is bounded from $L^2(\mathbb{R}^n)$ into $L^2(\mathbb{R}^n)$ with norm proportional to the best bound $C_{\alpha\beta}$.
    \label{theo:22boundR}
\end{theorem}

\begin{theorem}
    Let $T \in \Psi^m_{\rho, \delta}(\mathbb{T}^n \times \mathbb{Z}^n) $, $0 < \rho \leq 1$, $0 \leq \delta < 1$, with symbol $p(x, \xi)$ and with kernel

    \begin{equation}
        k(x, y) = \sum_{\xi \in \mathbb{Z}^n} e^{i2\pi(x - y) \cdot \xi} p(x, \xi).
    \end{equation}
    \begin{enumerate}
        \item[(a)] (Pseudo-local property) $k$ is smooth outside the diagonal. Moreover, given $\alpha, \beta \in \mathbb{N}^n_0$, then for any $N > (m + n + |\alpha + \beta|)/\rho$ we have
        \begin{equation}
            \sup_{x \neq y} |x - y|^N |\partial^\alpha_x \partial^\beta_y k(x, y)| = C_{\alpha\beta N} < \infty.
            \label{eq:teo3-1a}
        \end{equation}

        \item[(b)] Assuming  $p$ has compact support in $\xi$ uniformly in $x$, then $k$ is smooth and given $\alpha, \beta \in \mathbb{N}^n_0$, there is $C>0$ such that 
        \begin{equation}
            |\partial^\alpha_x \partial^\beta_y k(x, y)| \leq C\langle x - y \rangle^{-N}.
        \end{equation}


        \item[(c)] Assuming $m + M + n < 0$ for some $M \in \mathbb{Z}^+$, then $k$ is a bounded continuous function with bounded continuous derivatives up to order $M$.\\

        \item[(d)] Assuming $m + M + n = 0$ for some $M \in \mathbb{Z}^+$, then there is $ C >0$ such that 
        \begin{equation}
            \sup_{|\alpha + \beta| = M} |\partial^\alpha_x \partial^\beta_y k(x, y)| \leq C | \log |x-y| | \; , \quad x \neq y.
        \end{equation} 
    \end{enumerate}
\end{theorem}

\begin{theorem}
    Let $T \in \Psi^m_{\rho, \delta}(\mathbb{T}^n \times \mathbb{Z}^n) $, $0 < \rho \leq 1$, $0 \leq \delta < 1$, with symbol $p:=p(x, \xi)$ and with kernel $k:=k(x, y)$. Set $\lambda = \max\{ 0, (\delta - \rho)/2 \}$.
    \begin{itemize}
        \item[(a)] For any fixed $z \in \mathbb{T}^n$, and $\sigma \geq 1$, we have the kernel inequalities:
        \begin{equation}
            \sup_{|y-z| \leq \sigma} \int_{|x-z| > 2\sigma}|k(x, y) - k(x, z)|\diff x \leq C,
        \end{equation}
        \begin{equation}
            \sup_{|y-z| \leq \sigma} \int_{|x-z| > 2\sigma}|k(y, x) - k(z, x)|\diff x \leq C.
        \end{equation}
        \item[(b)] If $m \leq -n [(1-\rho)/2 + \lambda] $ and $\sigma < 1$, we have for any fixed $z \in \mathbb{T}^n$, 
        \begin{equation}
            \sup_{|y-z| \leq \sigma} \int_{|x-z| > 2\sigma^\rho}|k(x, y) - k(x, z)|\diff x \leq C.
        \end{equation}
        \item[(c)] If $m \leq -n (1 - \rho) / 2$ and $\sigma < 1$, we have  for any fixed $z \in \mathbb{T}^n$,
        \begin{equation}
            \sup_{|y-z| \leq \sigma} \int_{|x-z| > 2\sigma^\rho}|k(y, x) - k(z, x)|\diff x \leq C.
        \end{equation} 
    \end{itemize}
    \label{theo:sigma-kernel-estimate}
\end{theorem}

\begin{proof}
    \begin{itemize}
        \item[(a)] First, we notice that $|x-y| \geq |x - z| - |z - y| > \sigma$ in the domain of evaluation. Then by the triangle inequality and \cref{eq:teo3-1a} we have
        \begin{equation*}
            \sup_{|y-z| \leq \sigma} \int_{|x-z| > 2\sigma}|k(x, y) - k(x, z)|\diff x  
        \end{equation*}
        \begin{equation*}
            \leq \sup_{|y-z| \leq \sigma} \int_{|x-z| > 2\sigma}|k(x, y)|\diff x + \sup_{|y-z| \leq \sigma} \int_{|x-z| > 2\sigma}|k(x, z)|\diff x
        \end{equation*}
        \begin{equation*}
            \leq  \int_{|x-y|>\sigma} C_N|x-y|^{-N} \diff x + \int_{|x-z|>\sigma} C_N|x-z|^{-N} \diff x
        \end{equation*}
        \begin{equation*}
            \leq  \int_{\mathbb{T}^n} C_N\sigma^{-N} \diff x + \int_{\mathbb{T}^n} C_N\sigma^{-N} \diff x \leq C.
        \end{equation*}

        \item[(b)] As before, let $\tilde{p}$ be the corresponding symbol on $\mathbb{T}^n \times \mathbb{R}^n$.  Let $\varphi \in C_0^\infty(\mathbb{R})$ supported in $[1/2, 1]$ such that 
        \begin{equation*}
            \int_0^\infty\varphi(1/t)/t \diff t =\int_1^2\varphi(1/t)/t \diff t = 1.
        \end{equation*} 
        Define 
        \begin{equation*}
            k(x, y, t) = \int_{\mathbb{R}^n} e^{i2\pi(x - y) \cdot \xi} \tilde{p}(x, \xi) \varphi(\langle\xi\rangle/t) \diff \xi,
        \end{equation*}
        so that  
        \begin{equation*}
            k(x, y) = \int_0^\infty k(x, y, t) \diff t = \int_1^\infty k(x, y, t) \diff t.
        \end{equation*}
        Let $N > n/2$ be an integer, then we obtain the estimates
        \begin{equation*}
            \int_{|x-z|>2\sigma^\rho} |k(x, y, t) - k(x, z, t)| \diff x
        \end{equation*}
        \begin{equation*}
            \leq  \left[ \int_{\mathbb{T}^n} (1+t^{2\rho}|x-z|^2)^N |k(x, y, t) - k(x, z, t)| \diff x  \right]^{1/2} \left[ \int_{\mathbb{T}^n} (1+t^{2\rho}|x-z|^2)^{-N} \diff x \right]^{1/2} 
        \end{equation*}
        \begin{equation}
            \leq C\left[ \int_{\mathbb{T}^n} (1+t^{2\rho}|x-z|^2)^N |k(x, y, t) - k(x, z, t)| \diff x  \right]^{1/2}t^{-\rho n/2}.
            \label{eq:firstbound-b}
        \end{equation}
        So, for $|\alpha| \leq N$ we have that
        \begin{equation*}
            t^{\rho|\alpha|}(x-z)^\alpha \int_{\mathbb{R}^n} \left[ e^{i2\pi(x - y) \cdot \xi} - e^{i2\pi(x - z) \cdot \xi} \right]\tilde{p}(x, \xi)\varphi(\langle\xi\rangle/t)\diff \xi 
        \end{equation*}
        \begin{equation*}
            =t^{\rho|\alpha|}(x-z)^\alpha \int_{\mathbb{R}^n} e^{i2\pi(x - z) \cdot \xi} \left[ e^{i2\pi(z - y) \cdot \xi} - 1 \right]\tilde{p}(x, \xi)\varphi(\langle\xi\rangle/t)\diff \xi 
        \end{equation*}
        \begin{equation*}
            =\sum_{\beta\leq\alpha} C_{\alpha\beta} t^{\rho|\alpha|} \int_{\mathbb{R}^n} e^{i2\pi(x - z) \cdot \xi} \partial^\beta_\xi \left[ 
            \left(e^{i2\pi(z - y) \cdot \xi} - 1 \right) \tilde{p}(x, \xi) \right] \partial^{\alpha-\beta}_\xi \varphi(\langle\xi\rangle/t)\diff \xi.
        \end{equation*}
        Now, $|e^{i2\pi(x - z) \cdot \xi}-1| \leq |x-z||\xi|\leq t\sigma$ in the support of $\varphi(\langle\xi\rangle/t)$. On the other hand, $|\partial^\gamma_\xi e^{i2\pi(x - z) \cdot \xi}| \leq C_\gamma|y-z|^{|\gamma|} \leq C_\gamma \sigma ^{|\gamma|} \leq C\langle\xi\rangle^{-|\gamma|}(t\sigma)^{|\gamma|} $. Let us assume that $t\sigma < 1$, then for any $\chi \in C_0^\infty(\mathbb{R}^n)$ that is equal to $1$ on the support of $\varphi(\langle\xi\rangle/t)$ the set
        \begin{equation*}
            \Sigma_{\alpha\beta}=\left\{ \langle\xi\rangle^{n(1 - \rho)/2 + \rho|\beta|} \partial^\beta_\xi \left[\left( e^{i2\pi(x - z) \cdot \xi}-1 \right) 
            \tilde{p}(x, \xi) \right] \chi(\langle\xi\rangle/t) : |y-z| < \sigma,\; z \in \mathbb{T}^n \right\}
        \end{equation*}
        has measure bounded by $C\langle\xi\rangle^{n(1 - \rho)/2 + \rho|\beta|} (t\sigma)\langle\xi\rangle^{m-\rho |\beta|} \leq Ct\sigma \langle\xi\rangle^{-n\lambda} $. Thus, we can consider $\tilde{p}(x, \xi)$ as a symbol on $\mathbb{R}^n\times\mathbb{R}^n$ and by \cref{theo:22boundR} each of the corresponding operators with symbols $$\langle\xi\rangle^{n(1 - \rho)/2 + \rho|\beta|} \partial^\beta_\xi \left[\left( e^{i2\pi(x - z) \cdot \xi}-1 \right) 
            \tilde{p}(x, \xi) \right] \chi(\langle\xi\rangle/t),$$ on the set $\Sigma_{\alpha\beta}$ are bounded on $L^2(\mathbb{R}^n)$ with norm estimated by $Ct\sigma$.
        Hence \cref{eq:firstbound-b} can be estimated using Plancherel's identity by 
        \begin{equation}
            Ct\sigma\sum_{\beta\leq\alpha,\;|\alpha|\leq N} C_{\alpha\beta}t^{\rho|\alpha|}\left\|  
             \langle\xi\rangle^{-n(1-\rho)/2 - \rho|\beta|} t^{-|\alpha-\beta|} \partial^{\alpha-\beta}_\xi \varphi(\langle\xi\rangle/t) \right\|_{L^2(\mathbb{R}^n)}t^{-\rho n/2} 
        \end{equation}
        \begin{equation}
            \leq Ct^{\rho n/2}t^{-\rho n/2}t\sigma =Ct\sigma.
            \label{eq:tsigma<1-b}
        \end{equation}
        Now, let us dropp the restriction $t\sigma < 1$. For $|\alpha| = N$ we have that
        \begin{equation*}
            \int_{|x-z|>2\sigma^\rho} |k(x, y, t)| \diff x 
        \end{equation*}
        \begin{equation*}
            \leq \left[ \int_{\mathbb{T}^n} \left( t^{2\rho} |x-y|^2 \right)^N |k(x, y, t)|^2 \diff x \right]^{1/2} \left[ \int_{|x-y|>\sigma^\rho} \left(|t^{2\rho}|x-y|^2\right)^{-N}\diff x \right]^{1/2}
        \end{equation*}
        \begin{equation}
            \leq C \left[ \int_{\mathbb{T}^n} \left( t^{2\rho}|x-y|^2 \right)^N|k(x,y,t)|^2 \diff x \right]^{1/2}t^{-\rho N}\sigma^{\rho(n/2 - N)}.
            \label{eq:secondbound-b}
        \end{equation} 
        Let $|\alpha|=N$, then 
        \begin{equation*}
            t^{\rho|\alpha|} (x-y)^\alpha \int_{\mathbb{R}^n} e^{i2\pi(x - y) \cdot \xi} \tilde{p}(x, \xi)\varphi(\langle\xi\rangle/t)\diff\xi 
        \end{equation*}
        \begin{equation*}
            = \sum_{\beta\leq\alpha} C_{\alpha\beta}t^{\rho|\alpha|} \int_{\mathbb{R}^n} e^{i2\pi(x - y) \cdot \xi} \partial^\beta_\xi \tilde{p}(x, \xi)t^{-|\alpha-\beta|} \partial^{\alpha-\beta}_\xi \varphi(\langle\xi\rangle/t)\diff \xi,
        \end{equation*}
        and for each $\beta\leq\alpha$ the $L^2(\mathbb{T}^n)$ norm as a $x$-function of 
        \begin{equation*}
            \int_{\mathbb{R}^n} e^{i2\pi(x - y) \cdot \xi} \partial^\beta_\xi \tilde{p}(x, \xi)t^{-|\alpha-\beta|} \partial^{\alpha-\beta}_\xi \varphi(\langle\xi\rangle/t)\diff \xi
        \end{equation*}
        is equal to the $L^2(\mathbb{T}^n)$ norm as a $x$-function of 
        \begin{equation*}
            \int_{\mathbb{R}^n} e^{i2\pi x \cdot \xi} \partial^\beta_\xi \tilde{p}(x+y, \xi)t^{-|\alpha-\beta|} \partial^{\alpha-\beta}_\xi \varphi(\langle\xi\rangle/t)\diff \xi.
        \end{equation*}
        On the other hand, the set $\{ \langle\xi\rangle^{n(1-\rho)/2 + \rho|\beta|} \partial^\beta_\xi\tilde{p}(x + y, \xi) : y \in \mathbb{T}^n  \}$ has measure bounded by $C\langle\xi\rangle^{n(1-\rho)/2 + \rho|\beta|} \langle\xi\rangle^{m- \rho|\beta|} = C\langle\xi\rangle^{n(1-\rho)/2 +m}   $ so their respective operators are bounded on $L^2(\mathbb{R}^n)$ by \cref{theo:22boundR}. Thus, \cref{eq:secondbound-b} can be estimated by  

        \begin{equation*}
            C\sum_{\beta\leq\alpha,\; |\alpha|=N} t^{\rho|\alpha|}t^{-n(1-\rho)/2 - \rho|\beta|}t^{-\rho|\alpha-\beta|}t^{n/2}t^{\rho N} \sigma^{\rho(n/2 - N)} \leq C(t\sigma)^{\rho (n/2 - N)}.
        \end{equation*}
        In a similar way, we can estimate 
        \begin{equation*}
            \int_{|x-z|>2\sigma^\rho} |k(x, z, t)| \diff x \leq C(t\sigma)^{\rho(n/2 - N)}.
        \end{equation*}
        Using these estimates and \cref{eq:tsigma<1-b} we get the result from the following expression:
        \begin{equation}
            \int_{|x-z| > 2\sigma^\rho}|k(x, y) - k(x, z)|\diff x \leq C\left[ \int_1^{1/\sigma} t\sigma + \int_{1/\sigma}^\infty(t\sigma)^{\rho(n/2-N)} \right]\diff t/t \leq C
            \label{eq:lastbound-b}.
        \end{equation} 
        Thus, completing the proof.\\

        \item[(c)] First, let us notice that
        \begin{equation*}
            k(y, x, t) - k(z, x, t) 
        \end{equation*}
        \begin{equation*}
            =\int_{\mathbb{R}^n} e^{-i2\pi (x-y) \cdot \xi}[\tilde{p}(y, \xi) - \tilde{p}(z, \xi)]\varphi(\langle\xi\rangle/t) \diff \xi
        \end{equation*}
        \begin{equation*}
            +  
            \int_{\mathbb{R}^n} e^{-i2\pi x \cdot \xi} \left[ e^{i2\pi y \cdot \xi} -e^{i2\pi z \cdot \xi}  \right]\tilde{p}(z, \xi)\varphi(\langle\xi\rangle/t) \diff \xi
        \end{equation*}
        \begin{equation*}
            =f(x-y, y, z, t) + g(x, y, z, t).
        \end{equation*}
        Then, we obtain 
        \begin{align*}
            \int_{\mathbb{T}^n}|g(x, y, z, t)|\diff x \leq & Ct^{-\rho n/2} \left[ \int_{\mathbb{T}^n} |g(x, y, z, t)|^2(1 + t^{2\rho}|x|^2)^N \diff x  \right]^{1/2}\\
            \leq & Ct^{-\rho n/2} \sum_{|\alpha|\leq N}  \left[ \int_{\mathbb{T}^n} |(t^\rho x)^\alpha g(x, y, z, t)|^2\diff x  \right]^{1/2}.
        \end{align*}
        Let us notice that $g$ is the Fourier transform on $\mathbb{R}^n$ on the first variable of the function $G(\xi, y, z, t) = \left[ e^{i2\pi y \cdot \xi} -e^{i2\pi z \cdot \xi}  \right]\tilde{p}(z, \xi)\varphi(\langle\xi\rangle/t)$. Also, as above, $|\partial^\gamma_\xi( e^{i2\pi y \cdot \xi} -e^{i2\pi z \cdot \xi}  )| \leq Ct\sigma$ when $t\sigma<1$. Hence, assuming $t\sigma < 1$ and using Plancherel's identity, we get
        \begin{align*}
            \int_{\mathbb{T}^n}|g(x, y, z, t)|\diff x \leq &Ct^{-\rho n/2} \sum_{|\alpha|\leq N} \left\| \partial^\alpha_\xi G(\xi, y, z ,t) \right\|_{L^2(\mathbb{R}^n)}\\ 
            \leq &Ct^{-\rho n/2} \sum_{\beta\leq\alpha,\;|\alpha|\leq N} \left\| \partial^\beta\xi \left[ (e^{i2\pi y \cdot \xi} -e^{i2\pi z \cdot \xi})  \tilde{p}(z, \xi)\right]  \partial^{\alpha-\beta}_\xi\varphi(\langle\xi\rangle/t)  \right\|_{L^2(\mathbb{R}^n)} \\
            \leq & Ct^{-\rho n/2} \sum_{\beta\leq\alpha,\;|\alpha|\leq N} t\sigma t^{m-\rho|\beta|}t^{-|\alpha-\beta|}t^{n/2} \\
            \leq& Ct\sigma t^{n(1-\rho)/2 + m} \leq Ct\sigma.
        \end{align*}
        Now, we drop the restriction $t\sigma<1$ and we notice that
        \begin{equation*}
            g(x, y, z, t) = \int_{\mathbb{R}^n} e^{-i2\pi (x-z) \cdot \xi} \left[ e^{i2\pi (y-z) \cdot \xi} -1 \right]\tilde{p}(z, \xi)\varphi(\langle\xi\rangle/t) \diff \xi.
        \end{equation*}
        Thus, we have that
        \begin{equation*}
            \int_{\mathbb{T}^n} |g(x, y, z, t)|\diff x \leq Ct^{-\rho N}\sigma^{\rho(n/2 - N)} \left[ \int_{\mathbb{T}^n} (t^{2\rho}|x-z|^2)^N|g(x, y, z, t)|^2\diff x \right].
        \end{equation*}
        We can use Plancherel's identity as above, to get
        \begin{align*}
            \int_{\mathbb{T}^n} |g(x, y, z, t)|\diff x \leq & Ct^{-\rho N}\sigma^{\rho(n/2 - N)} \sum_{\beta\leq\alpha,\;|\alpha|=N} \left\| \partial^\beta_\xi \left[ (e^{i2\pi (y-z) \cdot \xi} -1)\tilde{p}(z, \xi) \right] \partial^{\alpha-\beta}_\xi\varphi(\langle\xi\rangle/t) \right\|_{L^2(\mathbb{R}^n)} \\
            \leq& Ct^{-\rho N}\sigma^{\rho(n/2 - N)} \sum_{\beta\leq\alpha,\;|\alpha|=N}t^{-|\alpha-\beta|}t^{n/2} \\
            \leq & C(t\sigma)^{\rho(n/2-N)}.
        \end{align*}
    \end{itemize}
    We can use the same procedures to find these bounds for $f(x-y, y,z, t)$ by noticing that $\int_{|x-z|>2\sigma^\rho} |f(x-y, y,z, t)|\diff x \leq \int_{|x|>\sigma^\rho} |f(x, y,z, t)|\diff x$. Then the desired bound comes from the calculation in \cref{eq:lastbound-b}.
\end{proof}

\end{document}
